\par\Needspace{6\baselineskip}\medskip
\begin{center}\hypertarget{native-222-article160}{}\pdfbookmark[2]{160.}{article-160}160.\end{center}
In setting out the method to which the preceding investigations were devoted, we encountered certain special cases in which it cannot be applied, at least in the form in which we presented it. We found that this difficulty arises:

\SourceLeaf{223}\noindent\emph{first}, whenever one of the three geocentric positions coincides either with the corresponding heliocentric position of the Earth or with the opposite point (the latter case can evidently occur only when the celestial body has passed between the Sun and the Earth),

\noindent\emph{second}, whenever the first geocentric position of the celestial body coincides with the third,

\noindent\emph{third}, whenever all three geocentric positions lie on the same great circle as the Earth's second heliocentric position.

In the first case the position of one of the great circles $AB$, $A'B'$, $A''B''$ will remain indeterminate. In the second and third cases the position of the point $B^*$ will remain indeterminate. The methods explained above for determining the heliocentric positions from the geocentric ones when the quantities $P$, $Q$ are regarded as known therefore cease to work in these cases. There is, however, an essential distinction. In the first case the difficulty is due to the method alone, while in the second and third it is due to the nature of the problem itself. In the first case the determination can certainly be made if the method is suitably modified. In the second and third it is altogether impossible, and the heliocentric positions will remain indeterminate. It will be worth examining these relations briefly. An exhaustive treatment would be of little use, especially since an exact determination of the orbit is impossible in all these special cases, where even the slightest observational errors would have an enormous effect on it. The same difficulty will still arise whenever the observations approximate one of these cases without corresponding to it exactly. This must therefore be borne in mind when selecting observations. Particular care must be taken to avoid any position in which the celestial body is both near a node and near opposition or conjunction. One must also avoid sets of observations in which, at the last observation, the celestial body has returned nearly to the same geocentric position as at the first, and finally those in which the great circle drawn from the Earth's middle heliocentric position to the body's middle geocentric position makes a very acute angle with the direction of geocentric motion and, as it were, grazes the first and third positions.
